\usepackage[utf8]{inputenc}
\usepackage[T1]{fontenc}
\usepackage[french]{babel}
\usepackage{fontspec}
\usepackage{iftex}
\usepackage{titlesec}
\usepackage{geometry}
\usepackage{lmodern}
\usepackage{setspace}
\usepackage{fancyhdr}
\usepackage{titling}
\usepackage{graphicx}
\usepackage{xcolor}
\usepackage{tikz}
\usetikzlibrary{tikzmark}
\usepackage{colortbl}
\usepackage[table]{xcolor}
\usepackage{makecell} % <— important

\usepackage{ifthen}
\newboolean{renvoiCorrection}




\makeatletter
\newdimen\Cline@savewidth
\newcommand{\Clinecol}[3]{% #1 = i-j, #2 = thickness (ex: 2pt), #3 = color (ex: white)
  % sauver réglages actuels
  \noalign{%
    \global\Cline@savewidth=\arrayrulewidth
    % sauver la couleur actuelle des traits (macro interne colortbl)
    \xdef\Cline@saved@arc{\CT@arc}%
    \global\arrayrulewidth=#2\relax
    \arrayrulecolor{#3}%
  }%
  % tracer la ligne partielle
  \cline{#1}%
  % restaurer réglages
  \noalign{%
    \expandafter\arrayrulecolor\expandafter{\Cline@saved@arc}%
    \global\arrayrulewidth=\Cline@savewidth
  }%
}
\makeatother

\usepackage{svg}
\usetikzlibrary{backgrounds}
\usetikzlibrary{calc}
\usepackage{setspace}
\usepackage{lastpage}
\usepackage{environ}

% Permet un ragged-right propre avec césures autorisées
\usepackage{ragged2e}
\AtBeginDocument{\RaggedRight}
\usepackage[none]{hyphenat}   % aucune césure
\usepackage{xurl}             % URLs jamais débordantes
\setlength{\emergencystretch}{3em} % sécurité anti-overfull

\usepackage[most]{tcolorbox} % enhanced + breakable

\usepackage{enumitem}
\usepackage{changepage} %% Change la geometrie de la page localement
\usepackage{mdframed}
\usepackage{xparse}

\usetikzlibrary{calc}
\usepackage{environ}

\renewcommand{\UrlFont}{\DINOTCond}


\newcommand{\includepicto}[1]{\raisebox{-1.5ex}{\includesvg[width=0.8cm]{#1}} }


\usepackage{colortbl}  % Nécessaire pour colorer les traits

\usepackage{needspace}
\usepackage{etoolbox}
\preto{\subsection}{\Needspace{5\baselineskip}} % Garanti que si subsection est trp bas (pas assez de place libre) on saut de page
\preto{\section}{\Needspace{5\baselineskip}} 

\usepackage{fancyhdr}
\pagestyle{fancy}

\fancyhf{} % On vide en-têtes et pieds


\usepackage{amsmath, amsfonts, amssymb, amsthm}

\geometry{a4paper, margin=2cm}
\pagestyle{empty}
\setlength{\parindent}{0pt}
\setstretch{1.1}

% Entête et pied de page
\pagestyle{fancy}
\fancyhf{}
\renewcommand{\headrulewidth}{0pt} % supprime la ligne dans le header

\makeatletter
\newcommand{\getauthor}{\@author}
\makeatother

% --- Faire apparaître les \subsection* et \subsubsection* dans le sommaire ---
\makeatletter
\let\oldsubsection\subsection
\renewcommand{\subsection}{%
  \@ifstar{\subsection@star}{\oldsubsection}% version étoilée / non étoilée
}
\newcommand{\subsection@star}[1]{%
  \oldsubsection*{#1}% garde exactement ta mise en forme visuelle
  \phantomsection
  \addcontentsline{toc}{subsection}{\protect\numberline{}#1}% entrée ToC sans numéro
}
\makeatother

\makeatletter
% --- Redéfinition subsubsection ---
\let\oldsubsubsection\subsubsection
\renewcommand{\subsubsection}{%
  \@ifstar{\subsubsection@star}{\oldsubsubsection}%
}
\newcommand{\subsubsection@star}[1]{%
  \oldsubsubsection*{#1}% garde le style visuel
  \phantomsection
  \addcontentsline{toc}{subsubsection}{\protect\numberline{}#1}%
}
\makeatother

% S'assurer que la profondeur de sommaire est suffisante
\setcounter{tocdepth}{3}